% !PRINTED 102
\weberpage{102}
%<<CONT
dem Falle $m = 8$ mindestens zwei verschiedene Primzahlen in $m$ enthalten sein.

Setzen wir:
\[
m = q_1\, q_2 \cdots q_\mu,
\]
wenn $m$ ungerade oder nur durch 4 theilbar ist,
\[
m = q_2\, q_3 \cdots q_\mu,
\]
wenn $m$ durch 8 theilbar ist, so dass $q_1 = 4$ sein kann und in der zweiten Formel $q_2 = 8$ ist, während die übrigen $q$ ungerade Primzahlen sind, so sind die Indexmoduln in den drei Fällen:
\[
\begin{aligned}
& q_1 - 1,\quad q_2 - 1,\ \ldots,\ q_\mu - 1\\
& 2\phantom{{}_1 -{}1,}\ ,\quad q_2 - 1,\ \ldots,\ q_\mu - 1\\
& 2,\quad 2\phantom{{}_1 -{}1}\ ,\quad q_3 - 1,\ \ldots,\ q_\mu - 1.
\end{aligned}
\]
Weil hiernach alle Indexmoduln durch 2 theilbar sind, so sind die Bedingungen für $m$ erschöpft.

Es ist [\S.\,21, (11)]:
\weq{1}{\delta_{k,h} = 2,\ i_{k,h} = 1,\ c_{k,h} = \tfrac{1}{2} c_h,\ k = 1, 2,\ h = 1, 2, \ldots \mu.}
Nun hat man die Grössen $e_{k,h}$ so zu bestimmen, dass aus der Matrix
\weq{2}{\begin{array}{cccc}
e_{1,1}, & e_{1,2}, & \ldots & e_{1,\mu}\\
e_{2,1}, & e_{2,2}, & \ldots & e_{2,\mu}
\end{array}}
wenigstens eine ungerade zweireihige Determinante gebildet werden kann.

Damit $\mathfrak{A}$ primärer Theiler von $\mathfrak{N}$ sei, kommt noch die Bedingung \S.\,21, 6. hinzu, wo nun zu unterscheiden ist, ob $m$ durch 8 theilbar ist oder nicht. Ist $m$ durch 8 theilbar, so ist die Anzahl der Indexmoduln um eins grösser, als die Anzahl der Primfactoren von $m$, und es giebt einen Indexmodul, dem kein Primfactor von $m$ entspricht; im anderen Falle stimmen beide Zahlen überein. Nach \S.\,21, 6. hat man die $e_{k,h}$ so zu wählen, dass in keinem der Paare
\[
e_{1,1}, e_{2,1};\ e_{1,2}, e_{2,2};\ \ldots\ e_{1,\mu}, e_{2,\mu},
\]
denen ein Primfactor von $m$ entspricht, beide Zahlen gerade sind. Wenn $m$ durch 8 theilbar ist, so ist ein dem Indexmodul $c_k = 2$ entsprechendes Paar darunter, dem keine solche Beschränkung auferlegt ist. Sind die Zahlen $e_{k,h}$ (nach dem Modul 2) so bestimmt, so
%CONT>>
